% ================================================================
% SECTION 4: GEOMETRY MODULE
% ================================================================
\section{Geometry Module}
\label{sec:geometry}

\begin{center}
\code{bionetflux.geometry.domain\_geometry}
\end{center}

The geometry module (1830 lines) manages network topology through three public types and four built-in geometry builders.

% ------------------------------------------------------------------
\subsection{Constants}

\begin{lstlisting}[language=Python, caption={Boundary type constants}]
EXTERIOR_BOUNDARY = -1   # Physical exterior boundary
PERIODIC_BOUNDARY = -2   # Periodic connections
SYMMETRY_BOUNDARY = -3   # Symmetry boundaries
\end{lstlisting}

These integer sentinels are used as the \code{domain2\_id} field in \code{ConnectionInfo} to flag boundary (non-interior) connections.

% ------------------------------------------------------------------
\subsection{ConnectionInfo Dataclass}

\begin{structurebox}
\begin{lstlisting}[language=Python]
@dataclass
class ConnectionInfo:
    domain1_id: int
    domain2_id: int          # < 0 for boundary connections
    parameter1: float        # parameter value in domain1
    parameter2: float        # parameter value in domain2
    metadata: Dict[str, Any] = None
\end{lstlisting}
\end{structurebox}

\noindent Query methods:
\begin{itemize}
  \item \code{is\_boundary\_connection() -> bool} --- \code{domain2\_id < 0}
  \item \code{is\_exterior\_boundary() -> bool}
  \item \code{is\_periodic\_boundary() -> bool}
  \item \code{is\_symmetry\_boundary() -> bool}
  \item \code{get\_boundary\_type() -> Optional[str]} --- returns \code{"exterior"}, \code{"periodic"}, \code{"symmetry"}, or \code{None}
\end{itemize}

% ------------------------------------------------------------------
\subsection{DomainInfo Dataclass}

\begin{structurebox}
\begin{lstlisting}[language=Python]
@dataclass
class DomainInfo:
    domain_id: int
    extrema_start: Tuple[float, float]   # (x1, y1) physical space
    extrema_end: Tuple[float, float]     # (x2, y2) physical space
    domain_start: float = 0.0            # parameter space start
    domain_length: float = 1.0           # parameter space length
    name: Optional[str] = None
    display_color: str = "blue"
    metadata: Dict[str, Any] = None
\end{lstlisting}
\end{structurebox}

\noindent If \code{domain\_length} is left at its default value of 1.0, the \code{\_\_post\_init\_\_} method recomputes it as the Euclidean distance between the extrema.

\smallskip\noindent Methods:
\begin{itemize}
  \item \code{euclidean\_length() -> float} --- $\sqrt{(x_2-x_1)^2+(y_2-y_1)^2}$
  \item \code{center\_point() -> Tuple[float, float]} --- midpoint of the extrema
  \item \code{direction\_vector() -> Tuple[float, float]} --- unit direction from start to end
\end{itemize}

% ------------------------------------------------------------------
\subsection{DomainGeometry Class}

\begin{lstlisting}[language=Python, caption={Constructor}]
class DomainGeometry:
    def __init__(self, name: str = "unnamed_geometry"):
        self.name = name
        self.domains: List[DomainInfo] = []
        self.connections: List[ConnectionInfo] = []
        self._next_id = 0
        self._global_metadata: Dict[str, Any] = {}
\end{lstlisting}

\subsubsection{Domain Management}

\begin{longtable}{p{8cm}p{7cm}}
\toprule
\textbf{Method} & \textbf{Description} \\
\midrule
\code{add\_domain(extrema\_start, extrema\_end, domain\_start=None, domain\_length=None, name=None, display\_color="blue", **metadata) -> int} & Add a domain; returns integer domain ID \\
\code{get\_domain(domain\_id) -> DomainInfo} & Retrieve domain by ID \\
\code{get\_all\_domains() -> List[DomainInfo]} & Get all domains \\
\code{num\_domains() -> int} & Number of domains \\
\code{find\_domain\_by\_name(name) -> Optional[int]} & Find domain ID by  name \\
\code{remove\_domain(domain\_id)} & Remove domain and its connections \\
\code{get\_domain\_names() -> List[str]} & List all domain names \\
\bottomrule
\end{longtable}

\subsubsection{Connection Management}

\begin{longtable}{p{8cm}p{7cm}}
\toprule
\textbf{Method} & \textbf{Description} \\
\midrule
\code{add\_connection(domain1\_id, domain2\_id, param1, param2, **meta) -> int} & Add inter-domain connection \\
\code{add\_exterior\_boundary(domain\_id, parameter, **meta) -> int} & Mark an exterior boundary \\
\code{add\_periodic\_boundary(domain\_id, parameter, **meta) -> int} & Mark a periodic boundary \\
\code{add\_symmetry\_boundary(domain\_id, parameter, **meta) -> int} & Mark a symmetry boundary \\
\code{get\_boundary\_connections() -> List} & All boundary connections \\
\code{get\_interior\_connections() -> List} & All interior connections \\
\code{get\_connection(id) -> ConnectionInfo} & Retrieve connection by ID \\
\code{remove\_connection(id)} & Remove connection \\
\code{get\_connections\_by\_type(type\_str) -> List} & Filter connections \\
\code{num\_connections() -> int} & Total connection count \\
\bottomrule
\end{longtable}

\subsubsection{Analysis and Validation}

\begin{longtable}{p{8cm}p{7cm}}
\toprule
\textbf{Method} & \textbf{Description} \\
\midrule
\code{validate\_geometry(verbose=False) -> bool} & Comprehensive validation: domain count, length positivity, extrema format, NaN/inf checks \\
\code{find\_intersections(tolerance=1e-6)} & Detect segment intersections \\
\code{get\_connectivity\_info(tolerance=1e-6) -> Dict} & Connectivity analysis \\
\code{suggest\_parameter\_spacing(gap=0.1)} & Suggest parameter space layout \\
\code{get\_bounding\_box() -> Dict[str, float]} & \code{\{x\_min, x\_max, y\_min, y\_max\}} \\
\code{total\_length() -> float} & Sum of all domain lengths \\
\code{summary() -> str} & Human-readable summary \\
\bottomrule
\end{longtable}

\noindent The class also implements \code{\_\_len\_\_}, \code{\_\_getitem\_\_}, and \code{\_\_iter\_\_} for Pythonic iteration.

\subsubsection{Self-Test}

\code{DomainGeometry} provides \code{create\_test\_geometries} (classmethod) and \code{run\_self\_test} for validation.

% ------------------------------------------------------------------
\subsection{Built-in Geometry Builders}
\label{sec:geometry_builders}

Four standalone functions create common topologies:

\begin{longtable}{p{7cm}p{8cm}}
\toprule
\textbf{Function} & \textbf{Description} \\
\midrule
\code{build\_arc\_sequence\_geometry(N=2, start=0.0, length=1.0)} & $N$ consecutive aligned segments with exterior boundaries at the two ends and trace-continuity connections at interior junctions \\
\code{build\_grid\_geometry(N=4)} & OoC-style grid: 4 vertical spines plus $4N$ horizontal connectors, with exterior boundaries and junction connections fully wired \\
\code{build\_labyrinth\_geometry(n\_cols=3, n\_rows=8, cell\_size=1.0)} & Serpentine maze structure \\
\code{create\_maze\_geometry(data\_dir=None, length=None)} & Load geometry from CSV files \code{points.csv} and \code{lines.csv}; each line in the CSV defines a domain between two named points.  Boundary points (those appearing in only one line) are marked as exterior boundaries.  The resulting boundary-point map is stored in the geometry's global metadata for use by \code{apply\_boundary\_overrides}. \\
\bottomrule
\end{longtable}
